\documentclass[11pt,a4paper]{article}
\usepackage[margin=1in]{geometry}
\usepackage{booktabs}
\usepackage{longtable}
\usepackage{graphicx}
\usepackage{hyperref}
\usepackage{listings}
\usepackage{xcolor}
\usepackage{enumitem}
\usepackage{caption}
\usepackage{amsmath}
\usepackage{amssymb}
\usepackage{tcolorbox}
\tcbuselibrary{breakable}

\hypersetup{
    colorlinks=true,
    linkcolor=blue,
    filecolor=magenta,
    urlcolor=cyan,
}

\tcbset{
    promptbox/.style={
        colback=blue!5!white,
        colframe=blue!75!black,
        title={#1},
        fonttitle=\bfseries,
        before skip=6pt,
        after skip=6pt,
        breakable,
    }
}

\title{Performance Report 22: Vulkan Engine \\
\large Vegetation, Frame Structure, and Memory Reservations --- the Subsystems Reports 17--21 Never Audited}
\author{}
\date{\today}

\begin{document}

\maketitle

\begin{abstract}
Reports~17--21 audited ray tracing, water frame machinery, water noise ALU,
and solid/water rasterization (twelve of report~21's fifteen items are now
implemented or soundly declined, see the status appendix there). Untouched
until now: the \emph{vegetation} subsystem, the \emph{frame structure} (18
async submits over $\sim$8 queues per frame), and the \emph{memory
reservations} (multi-GB fixed pools). This report audits those three areas
statically, with SPIR-V measurements and byte tallies from the creation code.

It finds \textbf{15 issues}. A fixed 128\,MB vegetation instance reservation
(4096 chunks $\times$ 2048 instances, always committed) sits beside $\sim$1\,GB
of solid mesh pools and $\sim$310\,MB of debug-only SDF/bounding-box streams,
all capacity-provisioned with no utilization feedback (C1). Vegetation
geometry is re-expanded up to 9 times per frame (depth, color, and three
shadow cascades $\times$ two paths), 24 vertices per instance with wind,
height noise and trigonometry evaluated per vertex instead of per instance
(C2, H4). The frame pays $\sim$18 queue submits and $\sim$9 capacity-sized
cull dispatches (C3). Behind them: an impostor fragment stage that writes
\texttt{gl\_FragDepth} and runs full CSM per pixel (H5), a provably exact
fetch-saving reorder in the foliage shader (H6), per-pixel CSM on grass that
belongs at vertex rate (H7), and a view-independent 2048$\times$1024 sky
equirect re-rendered every frame (H8). Each issue carries the standard
technique and a copy-pasteable prompt preserving validation-clean execution,
Synchronization2, and the AGENTS.md raster$+$CSM fallback.

Scope score: \textbf{6.0/10}; combined with report~21's 6.9/10 the application
stands at \textbf{6.4/10} (\S\ref{sec:score22}). The pattern of this engine
holds: submission architecture excellent, geometry processed too many times,
memory reserved too eagerly.
\end{abstract}

\tableofcontents
\newpage

% =====================================================================
\section{Executive Summary}
% =====================================================================

\subsection{Scope and Method}

Three areas, none previously reported on: (1) the vegetation subsystem
(\texttt{VegetationRenderer}, billboard/impostor shaders, impostor capture,
vegetation shadow cascades); (2) frame structure (submit count, cull
dispatch count, per-frame host churn, sky scheduling); (3) memory
reservations (pools, targets, debug streams, capture arrays). Ray tracing
(17--18), water machinery and noise (19--20), and solid/water rasterization
(21) are referenced only at interfaces. Method matches report~21: full reads
with file:line evidence, \texttt{spirv-dis} on shipped modules, byte tallies
from creation code. The application was not executed; dynamic magnitudes are
analytic estimates from verified structure.

\subsection{Measured Budgets}

\begin{table}[h]
\centering
\small
\begin{tabular}{@{}lrl@{}}
\toprule
\textbf{Module} & \textbf{Words} & \textbf{Role} \\
\midrule
\texttt{vegetation.vert.spv} & 4\,731 & billboard VS expansion (24 verts/instance) \\
\texttt{vegetation.geom.spv} & 6\,408 & \textbf{built, never bound} (dead, M9) \\
\texttt{vegetation.frag.spv} & 7\,066 & foliage shading + per-pixel CSM \\
\texttt{impostors.geom.spv} & 4\,663 & \textbf{built, never bound} (dead, M9) \\
\texttt{impostors.frag.spv} & 6\,931 & impostor shading + depth reconstruct \\
\texttt{vegetation\_shadow.vert.spv} & 3\,972 & shadow VS expansion \\
\texttt{sky.frag.spv} & 4\,108 & fullscreen sky, every frame \\
\bottomrule
\end{tabular}
\caption{Vegetation/sky module sizes. The two geometry-shader modules have
zero in-tree references.}
\end{table}

Per-frame fixed structure (Maximum-equivalent, shadows on): \textbf{18 queue
submits} (13 in \texttt{MyApp.cpp} + 5 in the shadow task) over $\sim$8
queues, and \textbf{$\sim$9 cull dispatches} over fixed capacities (solid
1536 + water 192 + brush pools + merged vegetation streams). Committed
reservations at 1080p (device + host, \S\ref{sec:vram}): \textbf{$\sim$1.7\,GB
+ $\sim$1.5\,GB}, dominated by fixed pools, not live geometry.

\subsection{Recommended Order}

Utilization feedback + pool rightsizing (C1) $\rightarrow$ vegetation
multiplicity and VS hoisting (C2, H4) $\rightarrow$ submit consolidation
(C3) $\rightarrow$ the exact foliage reorder (H6) $\rightarrow$ impostor/sky
scheduling (H5, H8) $\rightarrow$ dead-code deletion (M9). CPU items
(M10--M11) and hygiene (L13--L15) are independent.

% =====================================================================
\section{Cost Model}
% =====================================================================

\begin{equation}
C = \underbrace{\sum_{\mathrm{passes}} G_{\mathrm{passes}} \times V \times C_{\mathrm{vert}}}_{\text{\S3, reports 20--21}} + \underbrace{S \times C_{\mathrm{submit}} + D \times C_{\mathrm{dispatch}}}_{\text{\S4, this report}} + \underbrace{M_{\mathrm{res}} \times C_{\mathrm{pressure}}}_{\text{\S\ref{sec:vram}, this report}}
\end{equation}

with $S \approx 18$ submits, $D \approx 9$ capacity-sized cull dispatches,
and $M_{\mathrm{res}} \approx 3.2$\,GB committed regardless of scene size.
Reports~20--21 attacked the first term; this report attacks the second
(fixed per-frame CPU/GPU overhead) and prices the third.

% =====================================================================
\section{Memory Reservations}
\label{sec:vram}
% =====================================================================

Tallied from creation code at $1920\times1080$, 3 frames in flight (device
memory unless noted). Pools are \texttt{createBuffer} commitments, not
virtual reservations.

\begin{table}[h]
\centering
\small
\begin{tabular}{@{}llr@{}}
\toprule
\textbf{Resource} & \textbf{Sizing} & \textbf{Bytes} \\
\midrule
Solid vertex/index pools & $1536 \times (512+128)$\,KB & 983\,MB dev + 983\,MB host mirror \\
Vegetation instances & $4096 \times 2048 \times 16$\,B & 128\,MB dev \\
Water pools & $192 \times 640$\,KB & 123\,MB dev + host mirror \\
SDF debug streams & $500\mathrm{k} \times (20+64+20)$\,B $\times 3$ & 126\,MB host + 30\,MB dev \\
BBox debug streams & same as SDF & 126\,MB host + 30\,MB dev \\
Water targets (color/depth/body/column/back) & full-res $\times 3$ & $\sim$275\,MB dev \\
Solid + veg + brush + shadow targets & $\sim$8 full-res sets + CSM & $\sim$300\,MB dev \\
Sky equirect $2048\times1024$ & $\times 3$ & $\sim$25\,MB dev \\
Impostor capture arrays (post-capture) & $60 \times 256^2 \times (4+4+4)$\,B & 47\,MB dev \\
\midrule
Total & & $\sim$1.7\,GB dev + $\sim$1.5\,GB host \\
\bottomrule
\end{tabular}
\caption{Committed memory at 1080p. Debug streams alone exceed the entire
water target set; solid pools exceed everything else combined.}
\end{table}

% =====================================================================
\section{Critical Issues}
% =====================================================================

\subsection{C1: Multi-GB Fixed Reservations With No Utilization Feedback}
\label{sec:c1}

\textbf{Severity: Critical.}

Three independent capacity systems each commit worst-case memory at init and
never look back: solid pools at $1536 \times 640$\,KB $= 983$\,MB
(\texttt{SceneRenderer.cpp:66--76,1785--1788}), vegetation instances at
$4096 \times 2048 \times 16$\,B $= 128$\,MB
(\texttt{VegetationRenderer.cpp:preallocate}, called with max constants at
\texttt{SceneRenderer.cpp:1821}), and SDF/bounding-box debug streams at
$\sim$310\,MB mostly host-visible (\texttt{IndirectRenderer.cpp}, unconditional
straight-line allocation in \texttt{initSlots} from \texttt{MAX\_SDF\_CUBES} /
\texttt{MAX\_BBOX\_CUBES} $= 500\mathrm{k}$). Per-frame memoized counts exist
(\texttt{getMeshCount}, \texttt{getInstanceTotal}), but nothing compares
utilization against capacity: no warning, no downsize path, no telemetry in
any widget. A scene using 5\% of pool capacity costs the same as a full one,
and on shared-memory iGPUs the host-visible half ($\sim$1.5\,GB) is RAM gone.

\textit{Technique: measure, warn, tier.} Log utilization per pool at startup
and per minute; warn below 25\%; add a settings-tiered capacity ladder
(small/medium/large scenes) instead of one worst case; gate the debug
streams on their debug views (allocate on first enable, free on disable).

\begin{tcolorbox}[promptbox={Prompt --- C1: utilization feedback and tiered pools}]
Goal: committed memory tracks scene size. Steps: (1) log
used/capacity per pool after scene load + a $<25\%$ warning; (2) add
capacity tiers selected by scene stats; (3) allocate SDF/bbox streams
lazily on first debug-view enable. Files:
\texttt{vulkan/renderer/SceneRenderer.cpp},
\texttt{vulkan/renderer/VegetationRenderer.cpp},
\texttt{vulkan/renderer/IndirectRenderer.cpp}. Validation: VRAM/RAM deltas
at each tier, identical renders, validation-clean.
\\ \textbf{Implemented as:} (1) shipped --- \texttt{SceneRenderer::
logMemoryUtilization()} (mesh/vertex/index use vs capacity per pool,
vegetation chunks/instances, debug-stream state) runs after every scene
load/generation, warning below 25\% while non-empty; (3) shipped --- the
$\sim$310\,MB SDF/bbox streams allocate on first non-empty set via
\texttt{ensureSdfBboxBuffers()}, with dummy-bound descriptors,
conditional barriers (omitted barriers for never-written memory are
no-ops), and the pre-existing null-guarded consumers; (2) tiers declined
--- init-time capacities back hard-failure overflow semantics, so a wrong
tier costs holes or an assert, unverifiable without executing. The lazy
streams are the safe 310\,MB; tiers need a runnable sizing study first.
\end{tcolorbox}

\subsection{C2: Vegetation Geometry Is Re-Expanded up to 9 Times per Frame}
\label{sec:c2}

\textbf{Severity: Critical.}

Every vegetation draw re-runs the full VS expansion (24 vertices per
instance: wind, height noise, trigonometry, MVP): billboard depth
(\texttt{VegetationRenderer.cpp:drawDepth}), billboard color + impostor
color (\texttt{drawColor}), and three shadow cascades $\times$ billboard and
impostor shadow (\texttt{drawShadowCascade} per cascade). Default
(impostors off): 5 expansions; impostors on: 9. The shadow VS
(\texttt{vegetation\_shadow.vert.spv}, 3\,972 words) repeats the trig/height
work with wind early-outed --- the per-instance constants are recomputed
rather than reused across passes, and no pass caches anything.

\textit{Technique: expand once, reuse.} Depth-prepass-equivalent variings
already exist for the color pass; the structural fix is H4's per-instance
bake (compute once at generation, read 9 times). Short term: skip the shadow
re-expansion where the cascade depth precision exceeds billboard detail
(coarse outer cascades).

\begin{tcolorbox}[promptbox={Prompt --- C2: stop re-expanding vegetation per pass}]
Goal: one expansion per instance per frame. Steps: (1) implement H4's aux
per-instance buffer so every pass reads baked statics; (2) evaluate dropping
per-plane detail (6 $\rightarrow$ 3 planes) in the two outer shadow
cascades. Files: \texttt{vulkan/renderer/VegetationRenderer.cpp},
\texttt{shaders/vegetation.vert},
\texttt{shaders/vegetation\_shadow.vert}. Validation: shadow-vs-color A/B,
timestamp deltas on veg passes, validation-clean.
\end{tcolorbox}

\subsection{C3: $\sim$18 Queue Submits per Frame Across $\sim$8 Queues}
\label{sec:c3}

\textbf{Severity: Critical (CPU).}

\texttt{submitCommandBufferAsyncToQueue} appears 13 times in
\texttt{MyApp.cpp} (cull, brush-solid, sky, solid, vegetation $\times 2$,
SDF, bbox, water $\times 2$, brush-liquid, RT build) plus 5 inside the
shadow task (serial cull, 3 cascade rasterizations, blur/restore). Each
submit is a kernel call plus timeline-semaphore and validation overhead,
paid unconditionally every frame whether the payload is large or trivial
(SDF/bbox submits run with debug off; brush-liquid with no brush).

\textit{Technique: submit consolidation + empty-submit elision.} Skip submits
whose command buffer recorded nothing (SDF/bbox off, brush-liquid absent,
water-empty already has its fast path); merge the 3 cascade rasterizations
on single-queue hardware (the code already branches
\texttt{singleQueue} but still submits thrice); fold tiny serial submits
into neighbors where semaphore DAGs allow.

\begin{tcolorbox}[promptbox={Prompt --- C3: consolidate frame submits}]
Goal: fewer, fuller submits. Steps: (1) early-out empty debug/brush-liquid
submits; (2) single-queue cascade merge; (3) count submits/frame in the
overlay to prove the delta. Files: \texttt{MyApp.cpp},
\texttt{vulkan/renderer/ShadowRenderer.cpp}. Validation: submit counter
down, frame time down on CPU-bound scenes, validation-clean (semaphore DAG
unchanged in meaning).
\end{tcolorbox}

% =====================================================================
\section{High-Severity Issues}
% =====================================================================

\subsection{H4: VS Expansion Repeats Per-Instance Work $24\times$}
\label{sec:h4}

\textbf{Severity: High.}

The migration from geometry-shader to vertex-shader billboard expansion
(\texttt{vegetation.vert:3--5}) roughly doubled per-instance vertex ALU: the
GS evaluated wind 36 perlin2 per instance once, while the VS evaluates
\texttt{applyWindSkew} (3 perlin2 each, $\sim$60 ALU per \texttt{perlin2})
plus \texttt{vegetationHeightScale} plus \texttt{sin+cos(theta)} per corner
vertex --- $\sim$7k flops per instance, of which height, density hash and
trigonometry depend only on per-instance constants. The header comment
claiming height is ``computed once per instance'' is aspirational: VS
invocations cannot share it.

\textit{Technique: bake timeless statics at generation.} Height scale and
the density hash never change: append them to a small per-instance aux
buffer (per-chunk sized like the instance buffer, indexed identically).
Trigonometry goes to a 256-entry \texttt{sin}/\texttt{cos} LUT (static
reassignment of random orientations, invisible; needs A/B, not proof).

\begin{tcolorbox}[promptbox={Prompt --- H4: bake per-instance statics}]
Goal: per-instance work runs once per instance. Steps: (1) aux buffer with
height scale + density hash written by \texttt{processPendingChunks};
(2) VS reads instead of evaluating; (3) optional trig LUT. Files:
\texttt{vulkan/renderer/VegetationRenderer.cpp},
\texttt{shaders/vegetation.vert}, \texttt{shaders/impostors.vert}.
Validation: A/B identical (aux path is exact), VS word count down,
validation-clean.
\end{tcolorbox}

\subsection{H5: Impostor Fragment Stage Reconstructs Depth and Shades Full}
\label{sec:h5}

\textbf{Severity: High.}

\texttt{impostors.frag} pays, per pixel: color + normal + depth fetches, a
two-matvec depth reconstruction with two perspective divides
(\texttt{:62--73}), full Blinn-Phong with \texttt{pow}, and full CSM
\texttt{ShadowCalculationHard} --- on quads that then write
\texttt{gl\_FragDepth}, disabling early-$z$ for the whole pass. Two exact
reorders exist: the dither-fade discard (:79--85) is independent of the
depth work and can move before it (saves a fetch + 2 matvecs on every
dithered-out pixel); the depth-bounds discard (:64) already runs first.

\textit{Technique: discard early, shade once.} Hoist the dither test above
the depth block; evaluate vertex-rate CSM (the quad face normal is flat ---
shadow varies slowly across it) as follow-up.

\begin{tcolorbox}[promptbox={Prompt --- H5: cheapen the impostor fragment stage}]
Goal: no wasted reconstruction on discarded pixels. Steps: (1) move the
dither block above the depth block; (2) prototype flat-varying CSM.
Files: \texttt{shaders/impostors.frag}. Validation: transition-zone A/B
identical (discard order is immaterial), validation-clean.
\end{tcolorbox}

\subsection{H6: Foliage Discard Runs After Bandwidth It Could Skip}
\label{sec:h6}

\textbf{Severity: High.}

\texttt{vegetation.frag:62--93} fetches background albedo + normal
(\texttt{textureLod} $\times 2$) \emph{before} the \texttt{weight < 0.3}
discard, yet \texttt{weight = opacityWeight} $\times$ \texttt{mix(...)}
with the mix in $[0,1]$, so \texttt{opacityWeight < 0.3} implies discard
exactly. The existing \texttt{opacityWeight == 0.0} early-out (:66) proves
the pattern; extending the threshold to $0.3$ skips both background fetches
on every sparse pixel (grass silhouettes: the majority of foliage pixels)
with bit-identical results. Line 93 stays as the backstop for
\texttt{opacityWeight} $\geq 0.3$ pixels the background demotes.

\textit{Technique: exact early-out hoisting.} One-constant change with a
proof instead of a benchmark.

\begin{tcolorbox}[promptbox={Prompt --- H6: hoist the foliage discard threshold}]
Goal: two fewer fetches on sparse pixels, zero behavior change. Steps:
change :66 to \texttt{opacityWeight < 0.3}. Files:
\texttt{shaders/vegetation.frag}. Validation: pixel-identical screenshots
(proof above), validation-clean.
\end{tcolorbox}

\subsection{H7: Full CSM per Foliage Pixel Belongs at Vertex Rate}
\label{sec:h7}

\textbf{Severity: High.}

Both foliage shaders run the full cascaded shadow selection per pixel (up
to 2 matvecs + blend + fetches in \texttt{vegetation.frag:129--138} and
\texttt{impostors.frag:112--121}) for a signal that varies over meters
while the leaf detail varies over centimeters. The billboard VS already
computes per-vertex world position and plane normal; evaluating the shadow
there and interpolating the scalar costs $\sim$1/1000th on dense foliage
with no visible change (cascade blend zones are smooth by design).

\textit{Technique: vertex-rate shadow varying.} Compute
\texttt{ShadowCalculation(Hard)} per vertex (billboards) / per quad corner
(impostors), pass one float, multiply in the fragment stage.

\begin{tcolorbox}[promptbox={Prompt --- H7: vertex-rate foliage shadows}]
Goal: CSM cost off the per-pixel path. Steps: (1) shadow varying through
VS/GS stages; (2) keep the NdotL $> 0.01$ fast paths. Files:
\texttt{shaders/vegetation.vert}, \texttt{shaders/vegetation.frag},
\texttt{shaders/impostors.vert}, \texttt{shaders/impostors.frag}.
Validation: shadow-boundary A/B (cascade seams smooth), validation-clean.
\end{tcolorbox}

\subsection{H8: View-Independent Sky Equirect Re-Rendered Every Frame}
\label{sec:h8}

\textbf{Severity: High.}

\texttt{sky\_equirect.frag} documents its own output as ``view-independent''
--- it depends on sky parameters, sun and time only --- yet
\texttt{MyApp.cpp:1697} re-renders the fixed $2048\times1024$ target every
frame ($\sim$2M pixels of 4k-word sky shading) on the sky queue, plus the
fullscreen sky inside the solid pass. Camera motion never invalidates the
equirect; only sky settings, sun position and time-of-day do.

\textit{Technique: dirty-gated recompute.} Hash the sky parameter block per
frame; re-record the equirect pass only on change (plus one forced refresh
on resize). Static skies drop to zero per-frame equirect cost.

\begin{tcolorbox}[promptbox={Prompt --- H8: cache the sky equirect}]
Goal: equirect renders on change, not per frame. Steps: (1) parameter hash
in the sky task; (2) skip submit when clean; (3) keep the in-pass
fullscreen sky untouched. Files: \texttt{MyApp.cpp},
\texttt{vulkan/renderer/SkyRenderer.cpp}. Validation: time-lapse A/B
(sun sweeping still updates), validation-clean.
\end{tcolorbox}

% =====================================================================
\section{Medium-Severity Issues}
% =====================================================================

\subsection{M9: Dead Shaders, Pipelines, and Entry Points}
\label{sec:m9}

\textbf{Severity: Medium.}

Built every build, bound never: \texttt{vegetation.geom.spv} (6\,408 words),
\texttt{impostors.geom.spv} (4\,663) and
\texttt{impostors\_depth.geom.spv} (4\,265) have zero in-tree references
(the VS-expansion migration retired them); \texttt{perlin\_noise.comp.spv}
has no dispatcher; the sky sphere pipelines
(\texttt{skyPipeline}/\texttt{skyGridPipeline}, ``retained for
Solid360Renderer cubemap capture'') outlived Solid360Renderer itself
(removed; only comments mention it); \texttt{VegetationRenderer::render()}
and \texttt{SolidRenderer::render()} have no callers (the deferred
depth/color paths are the live ones). Cost is startup time, disk, and reader
confusion --- plus a real hazard: the dead \texttt{.geom} sources ($\sim$500
lines) describe an expansion strategy the shipped shaders contradict.

\textit{Technique: delete.} Remove sources, Makefile entries are automatic
(wildcards), and the unbound pipeline creations.

\begin{tcolorbox}[promptbox={Prompt --- M9: delete the dead weight}]
Goal: build only what binds. Steps: (1) delete the three \texttt{.geom}
sources + \texttt{perlin\_noise.comp} if unused; (2) remove sphere-pipeline
creation; (3) remove or document the dead \texttt{render()} entries.
Files: \texttt{shaders/}, \texttt{vulkan/renderer/SkyRenderer.cpp},
\texttt{vulkan/renderer/VegetationRenderer.cpp},
\texttt{vulkan/renderer/SolidRenderer.cpp}. Validation: identical build
minus the artifacts, validation-clean.
\end{tcolorbox}

\subsection{M10: Per-Chunk Staging Churn in Vegetation Streaming}
\label{sec:m10}

\textbf{Severity: Medium.}

Every streamed chunk update allocates \emph{two} buffers (host staging +
device instance buffer), submits one async copy, fences, and deferred-frees
(\texttt{VegetationRenderer.cpp:processPendingChunks}, up to 10 chunks per
frame = 20 VMA allocations + 20 frees + fence traffic). VMA takes a global
lock per allocation; steady-state streaming churns it continuously.

\textit{Technique: ring/batch staging.} One persistent staging ring sized to
the worst-case chunk, suballocated per upload; batch all chunk copies of a
frame into a single submit/fence.

\begin{tcolorbox}[promptbox={Prompt --- M10: batch vegetation uploads}]
Goal: zero allocations in the streaming path. Steps: (1) persistent
staging ring; (2) single submit per frame for all pending chunks. Files:
\texttt{vulkan/renderer/VegetationRenderer.cpp}. Validation: allocation
counters flat during streaming, identical instances, validation-clean.
\end{tcolorbox}

\subsection{M11: Per-Frame Host Churn Inventory}
\label{sec:m11}

\textbf{Severity: Medium.}

Small per-frame host writes, each trivial alone: shadow UBO staging
(\texttt{SHADOW\_CASCADE\_COUNT}+1 \texttt{UniformObject}s $\times$ frames,
plus one memcpy + copy per cascade per frame and descriptor copies reserved
at \texttt{frameCount}$\times 3 \times 14$), \texttt{updateWindParamsUBO}
($\sim$100\,B rewritten up to 5 times per frame: 3 cascade draws + depth +
color), and \texttt{writeVegChunkInfo} (a full chunk-table memcpy per
\texttt{prepareCull}). None stalls; together they are per-frame mapped-memory
traffic and API calls on the hot thread.

\textit{Technique: write-once / write-on-change.} UBO staging only when
values change (cascades are static per frame already -- hoist out of the
per-cascade loop); wind params once per frame; chunk table only after
consolidation (like \texttt{vegChunkInfoMap} already is).

\begin{tcolorbox}[promptbox={Prompt --- M11: still the per-frame host writes}]
Goal: host writes proportional to change, not frames. Steps: hoist the
three call families as described. Files:
\texttt{vulkan/renderer/ShadowRenderer.cpp},
\texttt{vulkan/renderer/VegetationRenderer.cpp}. Validation: identical
buffers (change-detection is exact), validation-clean.
\end{tcolorbox}

\subsection{M12: Vegetation Owns Full-Res Targets and a Composite Leg}
\label{sec:m12}

\textbf{Severity: Medium.}

Vegetation renders to its own full-resolution swapchain-format color + D32
depth ($\sim$50\,MB at 1080p) precisely so it can run async -- then pays the
deferred depth+color double geometry to fill them, plus the composite's
2$\times$2 depth taps and color mix per screen pixel. The targets are sized
for grass-blade detail the impostor path deliberately abandons at distance;
nothing tiers them (compare \texttt{waterRenderScale}, report 21 M10).

\textit{Technique: tier like water.} Add a vegetation render scale (default
1.0, dense presets lower), reusing the idle-recreate path; evaluate single
vs double geometry with the C2 prepass-gate analysis.

\begin{tcolorbox}[promptbox={Prompt --- M12: tier vegetation resolution}]
Goal: resolution proportional to need. Steps: settings scale +
idle-guarded recreate + composite already size-agnostic. Files:
\texttt{utils/Settings.hpp},
\texttt{vulkan/renderer/VegetationRenderer.cpp}. Validation: VRAM down,
grass-detail A/B, validation-clean.
\end{tcolorbox}

% =====================================================================
\section{Low-Severity / Hygiene Issues}
% =====================================================================

\subsection{L13: Per-Pixel Bayer Matrix Construction}
\label{sec:l13}

Both foliage fragment shaders build \texttt{const int M[16]} inside
\texttt{main()} --- 16 register initializers per pixel (twice: vegetation +
impostor). Hoist to file-scope \texttt{const}. Trivial, zero-risk.

\subsection{L14: Impostor Overlap Band Renders Twice}
\label{sec:l14}

Billboards render below $1.15\times$ impostor distance while impostors
render above $0.5\times$: the overlap band shades every pixel twice with
complementary dithering, and both discards kill early-$z$ there. Narrow the
band (0.85--1.15 is the historical shape; measure whether 0.95--1.05
suffices) or accept as the cross-fade price. Design tradeoff, documented
here so it stays a decision.

\subsection{L15: Redundant Wind UBO Writes and Stale Comments}
\label{sec:l15}

\texttt{updateWindParamsUBO} is identical work issued up to 5 times per
frame (fold into M11's hoist). Separately, the ``40\% of instances are empty
sentinel'' comment (\texttt{VegetationRenderer.hpp:59}) is stale --
generation skips empties, so no sentinel waste exists (verified); the
sentinel VS branch is dead defense. Update the comment, keep the branch.

% =====================================================================
\section{Implementation Roadmap and Expected Deltas}
% =====================================================================

\begin{table}[h]
\centering
\small
\begin{tabular}{@{}llc@{}}
\toprule
\textbf{Item} & \textbf{Change} & \textbf{Est.\ frame delta} \\
\midrule
C1 & Utilization feedback + tiers, lazy debug streams & $-0.3$--1\,GB committed \\
C2 & Veg bake + cascade detail tiers & $-50$--$80$\% veg vertex cost \\
C3 & Submit consolidation/elision & CPU-bound scenes improve \\
H4 & Per-instance bake (exact parts) & $-30$\% veg VS ALU \\
H5 & Impostor dither reorder + vertex CSM & $-1$ fetch, $-2$ matvecs on fade pixels \\
H6 & Foliage discard reorder (exact) & $-2$ fetches on sparse pixels \\
H7 & Vertex-rate foliage shadows & CSM off the foliage pixel path \\
H8 & Dirty-gated equirect & $-2$M sky px/frame when static \\
M9--M12 & Dead code, batching, host churn, veg scale & additive \\
L13--L15 & Hygiene & noise floor \\
\bottomrule
\end{tabular}
\end{table}

C1 is the only item that changes what the application \emph{costs to run}
anywhere (OOM headroom on small GPUs); C2/H4 change what grass costs per
frame; the rest is scheduling hygiene. None alters the lighting model, and
the exact items (H6, H8-cache-miss behavior, M11 change-detection) are safe
by construction.

% =====================================================================
\section{Overall Score}
\label{sec:score22}
% =====================================================================

\subsection{Rubric (Report-22 Scope)}

\begin{table}[h]
\centering
\small
\begin{tabular}{@{}lcp{7.2cm}@{}}
\toprule
\textbf{Category} & \textbf{Wt.} & \textbf{Score and justification} \\
\midrule
Vegetation efficiency & 30\% & \textbf{5.5} --- merged cull, indirect draws,
capture-on-demand and early discards are good; 24$\times$ VS ALU, 5--9
expansions/frame and per-pixel CSM drag it down. \\
Frame structure & 25\% & \textbf{7.0} --- explicit multi-queue DAG with
timeline semaphores and batched barriers; 18 submits and capacity-sized
redundant dispatches (restores) cost. \\
Memory provisioning & 30\% & \textbf{5.0} --- pools are leak-free and
race-free, but multi-GB worst-case commitments (including unconditional
debug streams) with no utilization feedback. \\
Startup \& dead weight & 15\% & \textbf{7.5} --- everything builds at init,
no runtime compilation; dead modules and pipelines linger. \\
\bottomrule
\end{tabular}
\end{table}

\subsection{Overall: 6.0/10 Scope, 6.4/10 Combined}

\begin{equation}
0.30(5.5) + 0.25(7.0) + 0.30(5.0) + 0.15(7.5) = \mathbf{6.0}
\end{equation}

Report~21 scored the raster core 6.9/10 ($\sim$45\% of frame cost by weight);
this scope covers the remaining $\sim$55\% at 6.0:

\begin{equation}
0.45(6.9) + 0.55(6.0) = \mathbf{6.4}
\end{equation}

The engine's verdict is unchanged in kind since report~21, extended in
scope: submission and lifetime architecture excellent, geometry processed
too many times, memory reserved too eagerly. Landing C1--C3 moves the
combined score toward $\sim$7.5/10 with no redesign.

\end{document}
